

\section{A larger state can preserve the information structure}
The distinction between a process and its information becomes especially clear
in the work of Beiglb\"ock, Pfl\"ugl and Schrott. Their information process
recursively augments today's value by the conditional law of tomorrow's
information process. Equal laws of suitable Markov lifts characterize their
notion of equivalence of filtered processes
\cite[Definitions~2.6--2.8, Proposition~3.6 and Theorem~3.9]{litcloseprobview}.
A lift may be large or measure-valued. Its existence does not imply that a
small sufficient state can be recovered from financial statements. Moreover,
transferring strategies requires attention to independent randomization and
the available stochastic basis \cite[Theorem~5.4 and Remark~5.7]{litcloseprobview}.
An implementer must decide whether an actor may randomize and which observations
are available, rather than deducing these permissions from equal outcome laws.
